\begin{proof}[Proof of \cref{obj:prop:published-class-inclusion}]
Fix \(0<s\leq1\). We use the published excess-variance identity
\citep[Proposition 2.3, equation (2.4), under Assumption 2.1 and Definition 2.2]{Cytrynbaum2026Limits}
and the expected absolute-supremum contraction inequality
\citep[Theorem 12(4), with Definition 2]{BartlettMendelson2002Complexities}
through \cref{obj:lem:published-prognostic-capacity,obj:thm:full-design-lower}.

\begin{enumerate}
\item For every integer \(d\geq2\), \cref{obj:lem:published-prognostic-capacity} supplies the proper inclusion
\[
\left\{
P:\ P\text{ is the Gaussian completion of some }
m\in\mathcal H_{d,s}\text{ specified in }\cref{obj:synth_4}
\right\}
\subsetneq\mathcal Q_{d,s}.
\]
This is the first assertion.
% lean: CausalSmith.Experimentation.BivariateSobolevDesignCapacity.published_prognostic_capacity

\item Fix integers \(n,d\geq2\) and a Borel assignment probability kernel \(\pi\) based on the original covariate array. The admissibility equivalence in
\cref{obj:lem:published-prognostic-capacity} states that
\(\pi\in\mathcal D_{n,d,s}\) precisely when
\[
\mathbb E_\pi[Z_i\mid X]=0
\quad\text{almost surely},\qquad i=1,\ldots,n,
\]
under independent uniform covariates, and
\[
Z\ \perp\!\!\!\perp\
\bigl((O_i(0),O_i(1))\bigr)_{i=1}^n
\ \bigm|\ X
\]
in the experiment formed from \(n\) independent copies of every single-unit law \(P\) having bounded strictly positive covariate density and finite potential-outcome second moments, followed by assignment from \(\pi(X)\). These are exactly the stated specialization of the published admissibility conditions.
% lean: CausalSmith.Experimentation.BivariateSobolevDesignCapacity.published_prognostic_capacity

\item Fix an integer \(n\geq2\). With the lower constant and comparison scale defined by
\[
c_s=\frac{2}{(48+32\pi^2)\,16384^s},
\qquad
b_s(n,d)=\min\left\{1,\left(\frac{\binom d2}{n}\right)^s\right\},
\]
\cref{obj:thm:full-design-lower} gives
\[
R_s(n,2)\geq c_s b_s(n,2).
\]
Since \(\binom22=1\), \(n\geq1\), and \(-s<0\),
\[
n^{-s}\leq1,
\qquad
b_s(n,2)=\min\{1,n^{-s}\}=n^{-s}.
\]
Substitution yields
\[
R_s(n,2)\geq
\frac{2}{(48+32\pi^2)\,16384^s}\,n^{-s}.
\]
% lean: CausalSmith.Experimentation.BivariateSobolevDesignCapacity.full_design_lower
% lean: CausalSmith.Experimentation.BivariateSobolevDesignCapacity.bScale_two

\item Fix integers \(n,d\geq2\) and \(m\in\mathcal H_{d,s}\). Define the constant assignment kernel on \(([0,1]^d)^n\) by
\[
\pi_{\mathrm{fair}}(\,\cdot\,)
=
\bigotimes_{i=1}^n
\left(\frac12\delta_{-1}+\frac12\delta_{+1}\right).
\]
Its signs are independent and fair, independently of the original covariate sample. Consequently, for \(i,i'\in\{1,\ldots,n\}\),
\[
\mathbb E_{\pi_{\mathrm{fair}}}[Z_iZ_{i'}\mid X]
=
\begin{cases}
1,&i=i',\\
0,&i\ne i',
\end{cases}
\]
and expansion of the finite square gives
\[
\mathbb E_{\pi_{\mathrm{fair}}}
\left[
\left(\sum_{i=1}^n Z_i m(X_i)\right)^2
\middle|X
\right]
=
\sum_{i=1}^n m(X_i)^2.
\]
Square-integrability of \(m\) justifies integrating this identity. Each covariate marginal is uniform, so
\[
\begin{aligned}
\frac4n\,\mathbb E
\left[
\left(\sum_{i=1}^n Z_i m(X_i)\right)^2
\right]
&=\frac4n\sum_{i=1}^n\mathbb E[m(X_i)^2]\\
&=4\int_{[0,1]^d}m(x)^2\,dx\\
&=4\mathbb E[m(X_1)^2].
\end{aligned}
\]
Finally, \cref{obj:lem:exact-reflection-budget} supplies
\[
\int_{[0,1]^d}|m(x)|^2\,dx\leq1.
\]
Since \(m\) is real-valued, multiplying this inequality by \(4\) proves the claimed upper bound.
% lean: CausalSmith.Experimentation.BivariateSobolevDesignCapacity.loss_fairSigns_eq
% lean: CausalSmith.Experimentation.BivariateSobolevDesignCapacity.sobolev_cube_second_moment_le_one
\end{enumerate}
\end{proof}
